\begin{name}
    {\tenchude}
    {\tendethi}
    {\tentruong}
    {\thoigian}
\end{name}
\Opensolutionfile{ansbook}[ans/ansbook-4]
\TN
\Opensolutionfile{ans}[ans/ans-4-T]
\begin{ex}%Câu 1
Cho cấp số cộng $\left(u_n\right)$ với $u_1=-1$ và $u_2=4.$ Giá trị của $u_3$ bằng
\choice
{\True $ 9$}
{$-16$}
{$ 7$}
{$-8$}
\loigiai{
Ta có: $ d=u_2-u_1=4-\left(-1\right)=5$ nên $u_3=u_2+d=4+5=9$.}
\end{ex}

\begin{ex}%Câu 2
Cho hàm số $ y=\dfrac{ax+b}{cx+d}$ với $ a,b,c \in\mathbb{R}$ có đồ thị như hình vẽ. Mệnh đề nào đúng?\\
\centerline{\begin{tikzpicture}[line join=round, line cap=round,>=stealth,thick]
        \tikzset{every node/.style={scale=0.9}}
        \draw[->] (-4.1,0)--(4.1,0) node[below left] {$x$};
        \draw[->] (0,-4.1)--(0,4.1) node[below left] {$y$};
        \draw (0,0) node [below left] {$O$}
        (1,0)node[below right]{$1$};
        \draw[dashed,thin] (1.01,-4)--(1.01,4);
        \begin{scope}
            \clip (-4,-4) rectangle (4,4);
            \draw[samples=200,domain=-4:0.99,smooth,variable=\x] plot (\x,{(1*(\x)+1)/(1*(\x)+-1)});
            \draw[samples=200,domain=1.01:4,smooth,variable=\x] plot (\x,{(1*(\x)+1)/(1*(\x)+-1)});
            \draw[dashed,thin] (-4,1/1)--(4,1/1);
        \end{scope}
\end{tikzpicture}}
\choice
{$y'>0,\forall x\ne 1$}
{$y'>0,\forall x\in\mathbb{R}$}
{$y'<0,\forall x\in\mathbb{R}$}
{\True $y'<0, \forall x\ne 1$}
\loigiai{
Từ đồ thị ta thấy $y'<0,\forall x\ne 1.$}
\end{ex}

\begin{ex}%Câu 3
Trong không gian $ Oxyz,$ cho điểm $ M\left(1;2;-1\right)$ và mặt phẳng $(P):x+2y+z=0.$ Mặt phẳng $(Q)$ qua $ M$ và song song với $(P)$ có phương trình là
\choice
{$ x+2y+z+4=0$}
{$ x+2y+z-1=0$}
{$ x+2y-z-6=0$}
{\True $ x+2y+z-4=0$}
\loigiai{
Mặt phẳng $(Q)$ có phương trình: $ x+2y+z+d=0$ với $ d\ne 0$\\
Điểm $ M\in(Q)$ nên suy ra $ 1+2.2+\left(-1\right)+d=0\Rightarrow d=-4$\\
Vậy phương trình mặt phẳng $(Q)$ là: $ x+2y+z-4=0.$}
\end{ex}

\begin{ex}%Câu 4
Số nghiệm nguyên của bất phương trình $\log_{0,5}\left(2x+6\right)\ge-5$ là
\choice
{\True $ 16$}
{$ 13$}
{$ 15$}
{$ 8$}
\loigiai{
Bất phương trình đã cho $\Leftrightarrow 0<2x+6\le{\left(0,5\right)^{-5}}\Leftrightarrow 0<2x+6\le 32\Leftrightarrow-3<x\le 13$\\
Vì $ x\in\mathbb{Z}$ nên có tất cả $ 16$ nguyện nguyên thỏa mãn.}
\end{ex}

\begin{ex}%Câu 5
Cho hình lập phương $ABCD.A'{B}'{C}'{D}'$ có cạnh $ 2$ (tham khảo hình vẽ bên dưới). Độ dài của vectơ $\vec{u}=\vec{A'{C}'}-\vec{A'A}$ bằng\\
\centerline{
\begin{tikzpicture}[line cap=round,line join=round, >=stealth,scale=.7]
    \def \a{-1.5} \def \b{-1}\def \c{4.5} \def \h{4}
    \path (.5,.5)coordinate(A) 
    +(\a,\b)coordinate(B)
    +(\c,0)coordinate(D)
    ($(B)+(D)-(A)$)coordinate(C)
    +(0,\h) coordinate(C')
    ($(B)+(C')-(C)$)coordinate(B')
    ($(A)+(C')-(C)$)coordinate(A')
    ($(D)+(C')-(C)$)coordinate(D');
    \draw [dashed] (A)--(B)(D)--(A)--(A');
    \draw (B')--(B)--(C)(B')--(C')--(C)--(D)--(D')--(A')--(B')(C')--(D');
    \foreach \x/\g in {A/135,B/-135,C/-45,D/0,A'/135,B'/180,C'/-20,D'/0}\fill[red] (\x) circle (1pt)+(\g:3mm) node[black]{$\x$};
\end{tikzpicture}
}
\choice
{$ 2\sqrt{2}$}
{$\sqrt{3}$}
{$ 2\sqrt{6}$}
{\True $ 2\sqrt{3}$}
\loigiai{
Ta có $\vec{u'}=\vec{A{C}'}-\vec{A'A}=\vec{A{A}'}+\vec{A'{C}'}=\vec{A{C}'}\Rightarrow\left|\vec{u'}\right|=\left|\vec{A{C}'}\right|=A{C}'$\\
Khi đó $A{C}'=\sqrt{A{A^2}+A'{C^{\prime 2}}}=\sqrt{2^2+8}=2\sqrt{3}$}
\end{ex}

\begin{ex}%Câu 6
Biết $\displaystyle\int{f(x)\mathrm{\,d}x=\cos x+C}$ thì $\displaystyle\int{f'(x)\mathrm{\,d}x}$ bằng
\choice
{$\sin x+C'$}
{$\cos x+C'$}
{\True $-\sin x+C'$}
{$-\cos x+C$}
\loigiai{
Ta có $ f(x)=\left(\displaystyle\int{f(x)\mathrm{\,d}x}\right)'=\left(\cos x+C\right)'=-\sin x$ nên $\displaystyle\int{f(x)\mathrm{\,d}x=-\sin x+C}$}
\end{ex}

\begin{ex}%Câu 7
Cho hàm số $f(x)$ xác định trên $\left(-\infty ;0\right)\setminus\left\{-2\right\}$ và có bảng biến thiên bên dưới. Đồ thị hàm số đã cho có tổng số đường tiệm cận ngang và tiệm cận đứng là\\
\centerline{
\begin{tikzpicture}[>=stealth]
    \tkzTabInit[espcl=3,lgt=1.5,nocadre=false]
    {$x$/0.7,$f'(x)$/1,$f(x)$/3}
    {$-\infty$,$-2$,$+\infty$}
    \tkzTabLine{,+,d,-,}
    \tkzTabVar{-/$-\infty$,+D+/$+\infty$/$1$,-/$0$}
\end{tikzpicture}
}
\choice
{\True $ 1$}
{$ 3$}
{$ 0$}
{$ 2$}
\loigiai{
Từ đồ thị hàm số không có tiệm cận ngang và có một tiệm cận đứng là đường thẳng $ x=-2$.}
\end{ex}

\begin{ex}%Câu 8
Khi cắt vật thể bởi mặt phẳng vuông góc với trục $ Ox$ tại điểm có hoành độ là $ x$ $\left(0\le x\le 3\right),$ ta được mặt cắt là một hình vuông có cạnh là $\sqrt{9-x^2}$ (được mô hình hóa bởi hình vẽ bên dưới). Thể tích của vật thể đó bằng\\
\centerline{\includegraphics{images/4.8}}
\choice
{$ 171\pi $}
{$ 171$}
{$ 18\pi $}
{\True $18$}
\loigiai{
Thể tích vật thể là: $V=\displaystyle\int\limits_0^3S(x)\mathrm{\,d}x=\displaystyle\int\limits_0^3\left(9-x^2\right)\mathrm{\,d}x=18$ .}
\end{ex}

\begin{ex}%Câu 9
Cho lăng trụ tam giác $ ABC.A'{B}'{C}'.$ Biết diện tích mặt bên $ AB{B}'{A}'$ bằng $ 15$ và khoảng cách từ $ C$ đến mặt phẳng $\left(AB{B}'{A}'\right)$ bằng $ 6$ (tham khảo hình vẽ bên cạnh). Thể tích của khối lăng trụ $ ABC.A'{B}'{C}'$ bằng bao nhiêu?\\
\centerline{
    \begin{tikzpicture}[scale=1, font=\footnotesize, line join=round, line cap=round, >=stealth]
        \def\ac{4} % cạnh AC
        \def\ab{2} % cạnh AB
        \def\h{4} % chiều cao
        \def\gocA{50} % góc A của đáy
        \path
        (0,0) coordinate (A)
        (\ac,0) coordinate (C)
        (-\gocA:\ab) coordinate (B)
        ($(A)+(80:\h)$) coordinate (A')
        ($(B)-(A)+(A')$) coordinate (B')
        ($(C)-(A)+(A')$) coordinate (C');
        \draw (A')--(A)--(B)--(C)--(C')--(A')--(B')--(C') (B)--(B');
        \draw[dashed] (A)--(C);
        \foreach \x/\g in {A/180,B/-90,C/0,A'/180,C'/0,B'/-40}\fill[red] (\x) circle (1pt)+(\g:3mm) node[black]{$ \x $};
    \end{tikzpicture}
}
\choice
{$60$}
{\True $45$}
{$90$}
{$30$}
\loigiai{
Ta có: $V_{C.AB{B}'{A}'}=\dfrac{1}{3}.S_{AB{B}'{A}'}.d\left(C,\left(AB{B}'{A}'\right)\right)=\dfrac{1}{3}.15.6=30$\\
Mặt khác: $V_{C.A'{B}'{C}'}=\dfrac{1}{3}.d\left(C,\left(A'{B}'{C}'\right)\right).S_{A'{B}'{C}'}=\dfrac{1}{3}$ thể tích khối lăng trụ\\
Suy ra $V_{C.AB{B}'{A}'}=\dfrac{2}{3}$ thể tích khối lăng trụ nên $V_{ABC.A'{B}'{C}'}=\dfrac{3}{2}.30=45$.}
\end{ex}

\begin{ex}%Câu 10
Trong không gian $Oxyz,$ cho hai điểm $A\left(1;3;2\right)$ và $B\left(4;5;6\right).$ Gọi $\alpha $ là góc giữa đường thẳng $AB$ và mặt phẳng $\left(Oxy\right).$ Giá trị của $\text{cos}\alpha $ bằng
\choice
{$\dfrac{4\sqrt{29}}{29}\cdot $}
{$\dfrac{16}{29}\cdot $}
{\True $\dfrac{\sqrt{377}}{29}\cdot $}
{$\dfrac{13}{29}\cdot $}
\loigiai{
Vectơ pháp tuyến của mặt phẳng $\left(Oxy\right)$ là: $\vec{n_{\left(Oxy\right)}}=\left(0;\,0;\,1\right)$\\
Vectơ chỉ phương của đường thẳng $ AB$ là: $\vec{n_{AB}}=\left(3;\,2;\,4\right)$\\
Ta có: $\sin\alpha=\dfrac{\left|\vec{n}.\vec{u}\right|}{\left|\vec{n}\right|.\left|\vec{u}\right|}=\dfrac{4}{1.\sqrt{9+4+16}}=\dfrac{4}{\sqrt{29}}$ suy ra $\cos\alpha=\sqrt{1-\sin^2\alpha}=\dfrac{\sqrt{377}}{29}$.}
\end{ex}

\begin{ex}%Câu 11
Theo thống kê điểm trung bình môn Toán của một số học sinh đã trúng tuyển vào lớp 10 năm học 2024 – 2025 của một trường được kết quả như bảng sau:
\begin{center}
\begin{tabular}{|c|c|c|c|c|c|c|c|}
    \hline
    Khoảng điểm & $[6,5; 7)$ & $[7; 7,5)$ & $[7,5; 8)$ & $[8; 8,5)$ & $[8,5; 9)$ & $[9; 9,5)$ & $[9,5; 10)$ \\
    \hline
    Tần số & 7 & 10 & 17 & 24 & 13 & 8 & 5 \\
    \hline
\end{tabular}
\end{center}

Khoảng tứ phân vị của mẫu số liệu ghép nhóm trên là (Kết quả làm tròn đến hàng phần chục)
\choice
{\True $\Delta_Q=1,1$}
{$\Delta_Q=1$}
{$\Delta_Q=1,2$}
{$\Delta_Q=0,6$}
\loigiai{
Ta có: $ n=84$ và $\Delta Q=Q_3-Q_1$.\\
Mặt khác: $\dfrac{3n}{4}=63\Rightarrow{Q_3}=\dfrac{\dfrac{3n}{4}-\left(7+10+17+24\right)}{13}.0,5+8,5=\dfrac{113}{13}$\\
$\dfrac{n}{4}=21\Rightarrow{Q_1}=\dfrac{\dfrac{n}{4}-7-10}{17}.\left(8-7,5\right)+7,5=\dfrac{259}{34}$\\
Vậy khoảng tứ phân vị là: $\Delta Q=Q_3-Q_1=\dfrac{113}{13}-\dfrac{259}{34}\approx 1,1$}
\end{ex}

\begin{ex}%Câu 12
Một người gửi tiết kiệm $ 10$ triệu đồng vào một ngân hàng với lãi suất $7\%/$ một năm. Biết rằng nếu không rút tiền ra khỏi ngân hàng thì cứ sau mỗi năm, số tiền lãi sẽ được nhập vào vốn ban đầu. Sau $5$ năm mới rút lãi thì người đó thu được số tiền lãi là
\choice
{$14,026$ triệu đồng}
{$50,7$ triệu đồng}
{\True $4,026$ triệu đồng}
{$3,5$ triệu đồng}
\loigiai{
Số tiền cả gốc và lãi sau $ 5$ năm là: $P_5=P_0\left(1+r\right)^5=10\left(1+7\%\right)^5\approx 14,026$ triệu đồng.\\
Vậy số tiền lãi mà người đó nhận là: $ 14,026-10=4,026$ triệu đồng}
\end{ex}
\Closesolutionfile{ans}
\TNTF
\Opensolutionfile{ans}[ans/ans-4-TF]
\begin{ex}%Câu 13
Cho hàm số $y=\dfrac{a{x^2}+bx+c}{mx+n}$ có đồ thị như hình vẽ sau:\\
\centerline{
    \begin{tikzpicture}[line join=round, line cap=round,>=stealth,thick]
        \tikzset{every node/.style={scale=0.9}}
        \draw[->] (-4.1,0)--(4.1,0) node[below left] {$x$};
        \draw[->] (0,-4.1)--(0,4.1) node[below left] {$y$};
        \draw (0,0) node [below left] {$O$};
        \foreach \x/\nx in {-2/-2,-1/-1}
        \draw[thin] (\x,1pt)--(\x,-1pt) node [below] {$\nx$};
        \foreach \y/\ny in {-2/-2,1/1}
        \draw[thin] (1pt,\y)--(-1pt,\y) node [right] {$\ny$};
        \draw[dashed,thin](-2,0)--(-2,-2)--(0,-2);
        \draw[dashed,thin] (-0.99,-4)--(-0.99,4);
        \begin{scope}
            \clip (-4,-4) rectangle (4,4);
            \draw[samples=200,domain=-4:-1.01,smooth,variable=\x] plot (\x,{(1*((\x)^2)+2*(\x)+2)/(1*(\x)+1)});
            \draw[samples=200,domain=-0.99:4,smooth,variable=\x] plot (\x,{(1*((\x)^2)+2*(\x)+2)/(1*(\x)+1)});
            \draw[dashed,thin] (-4.1,-3.1)--(4.1,5.1);
        \end{scope}
    \end{tikzpicture}
}
\choiceTFt
{Hàm số đã cho nghịch biến trên khoảng $(-2;0)$}
{\True Đồ thị của hàm số đã cho có tiệm cận xiên $y=x+1$}
{Gọi $A,B$ là hai điểm cực trị của hàm số đã cho, diện tích của tam giác $OAB$ bằng $8$ (với $O$ là gốc tọa độ)}
{\True Một trục đối xứng của đồ thị đã cho là $d:y=\left(x+1\right)\tan\dfrac{3\pi}{8}\cdot $}
\loigiai{
a) Sai: Đồ thị đi xuống trên các khoảng $\left(-2;-1\right),\left(-1;0\right)$ nên nghịch biến trên các khoảng này.\\
b) Đúng: Tiệm cận xiên qua hai điểm $\left(-1;0\right),\left(0;1\right)$ là $ y=x+1$.\\
c) Sai: Hai điểm cực trị của đồ thị là $ A\left(0;2\right),B\left(-2;-2\right)\Rightarrow{S_{OAB}}=\dfrac{1}{2}\left| 0.\left(-2\right)-\left(-2\right).2\right|=2$.\\
d) Đúng: Tiệm cận đứng $ x=-1$. Trục đối xứng của đồ thị hàm số là đường phân giác của góc tạo bởi tiệm cận đứng và tiệm cận xiên:\\
$\dfrac{x-y+1}{\sqrt{2}}=\pm\left(x+1\right)\Leftrightarrow y=\left(x+1\right)\pm\sqrt{2}\left(x+1\right)=\left(1\pm\sqrt{2}\right)\left(x+1\right)=\tan\left(\dfrac{3\pi}{8}\right)$.}
\end{ex}
\begin{ex}%[Nguyễn Tuấn, dự án sáng tác đề 12 theo chủ đề]%[2H5V2-7] 
    Trong không gian với hệ tọa độ $O x y z$, một cabin cáp treo xuất phát từ điểm $A(10 ; 3 ; 0)$ và chuyển động đều theo đường cáp có véc-tơ chỉ phương là $\vec{u}=(2 ;-2 ; 1)$ với tốc độ là 4,5 m/s (đơn vị trên mỗi trục tọa độ là mét) (Hình bên dưới).
    \definecolor{ecru}{rgb}{0.76, 0.7, 0.5}
    \definecolor{darkolivegreen}{rgb}{0.33, 0.42, 0.18}
    \definecolor{deepskyblue}{rgb}{0.0, 0.75, 1.0}
    \definecolor{antiquebrass}{rgb}{0.8, 0.58, 0.46}
    \definecolor{arsenic}{rgb}{0.23, 0.27, 0.29}
    \definecolor{ashgrey}{rgb}{0.7, 0.75, 0.71}\definecolor{alizarin}{rgb}{0.82, 0.1, 0.26}
    \begin{center}
        \begin{tikzpicture}[line join=round, line cap=round,scale=1,transform shape]
            \clip (-3,-3) rectangle (3,3);
            \fill[bottom color=white,top color=deepskyblue!90, middle color=white] (-3,-3) rectangle (3,3);
            \tikzset{dat/.pic={
                    \def\N{ 
                        (-3,3)
                        ..controls +(20:.7) and +(150:.7) ..(-.8,1.5)
                        ..controls +(-30:.7) and +(150:.6) ..(1.5,-.8)
                        ..controls +(-30:1) and +(90:.4) ..(3,-3)--(-3,-3)--cycle
                        ;
                    }
                    \draw[black]\N;
                    \fill[darkolivegreen!50] \N;
            }}
            \tikzset{cap_treo/.pic={
                    \def\X{ 
                        (-.85,1)--(-.8,1)
                        ..controls +(-90:.9) and +(-180:.6) ..(-.1,0.1)--(-.1,.05)
                        ..controls +(-180:.65) and +(-90:.9) ..cycle
                        (-.15,0.05)--(-.15,-.1)--(-.23,-.22)--(-.6,-.22)--(-.5,-.1)--(-.5,.08)
                        ;
                    }
                    %\fill[black] \X;
                    \draw\X;
                    \draw(-3,2.3)--(3.5,-1.2);
                    \draw(-1.2,1)--(-.5,1);
                    \draw(-.5,1)--(-.9,1.2)--(-.95,1.17)--(-.6,.97)--cycle;
                    \draw[fill=arsenic!80](-.9,-.05)--(.52,-.05)--(.28,-.3)--(-1.05,-.3)--cycle;
                    \def\M{ 
                        (-.85,1)--(-.8,1)
                        ..controls +(-90:.9) and +(-180:.6) ..(-.1,0.1)--(-.1,.05)
                        ..controls +(-180:.65) and +(-90:.9) ..cycle
                        (-.15,0.05)--(-.15,-.1)--(-.23,-.22)--(-.6,-.22)--(-.5,-.1)--(-.5,.08)
                        ;
                    }
                    \fill[arsenic] \M;
                    \draw\M;
                    \def\N{ 
                        (-1.05,-.3)
                        ..controls +(-120:.6) and +(120:.6) ..(-.98,-1.5)--(.4,-1.5)
                        ..controls +(70:.6) and +(-60:.3) ..(.28,-.3)--cycle
                        ;
                    }
                    \fill[arsenic] \N;
                    \draw\N;
                    \def\P{ 
                        (.4,-1.5)
                        ..controls +(70:.6) and +(-60:.3) ..(.28,-.3)--(.52,-.05)
                        ..controls +(-40:.4) and +(50:.4) ..(.6,-1.2)--cycle
                        ;
                    }
                    \fill[arsenic!80] \P;
                    \draw\P;
                    \draw (-1.23,-1)--(.5,-1)--(.77,-.6);
                    \draw (-1.05,-.5)--(-1,-.4)--(.2,-.4)--(.25,-.5)--cycle
                    (.25,-.5)--(.3,-.6)--(.3,-.6)--(-1.1,-.6)--(.-1.05,-.5)--cycle
                    ;
                    \draw[alizarin](.42,-1)..controls +(100:.4) and +(-65:.4) ..(.2,-.3);
                    \draw[alizarin](-1.15,-1)..controls +(100:.3) and +(-120:.3) ..(-.95,-.3);
            }}
            \path
            (0,0)pic[scale=1]{dat}(0,0)pic[scale=1]{cap_treo};
        \end{tikzpicture}
        \begin{tikzpicture}[scale=1, font=\footnotesize, line join=round,xscale=.2, line cap=round,>=stealth]
            \def\a{1/16} 
            \def\xmin{-3} \def\xmax{12}
            \def\ymin{-3} \def\ymax{3} 
            \coordinate (O) at (0,0);
            \coordinate (E) at (-10,-3);
            \coordinate (N) at ($(E)!.7!(O)$);
            \coordinate (P) at ($(E)!.2!(O)$);
            \coordinate (D) at ($(E)!.3!(O)$);
            \coordinate (A) at ($(N)+(3,0)$);
            \coordinate (B) at (-14,1);
            \coordinate (M) at ($(A)!.7!(B)$);
            \draw[->] (\xmin,0)--(\xmax,0) node [above right]{$y$};
            \draw[->] (O)--(0,\ymax) node [left]{$z$};
            \draw[->] (O)--(E) node [below right]{$x$};
            \node at (0,0)[above right]{$O$};
            \draw[dashed] (B)--(P)node [right]{$550$} (N)--(A)node [below right]{$A(10;3;0)$}--(3,0);
            \draw(B)node [above]{$B$}--(A);
            \draw[->,red](O)--(-4,.6)node [above right]{$\vec{u}$};
            \node at (N) [left]{$10$};
            \node at (M) [above]{$M$};
            \node at (P) [left]{$x_B$};
            \draw[fill=black] (3,0) circle (.5pt) node[below]{\footnotesize $3$};
            %	\path (-28,0) node[opacity=.5,scale=.5] {\includegraphics{images/h35}};
            \clip (\xmin+0.1,\ymin+0.1) rectangle (\xmax-0.1,\ymax-0.1);
        \end{tikzpicture}
    \end{center}
    \choiceTF
    {Phương trình chính tắc của đường cáp là $\dfrac{x+10}{2}=\dfrac{y+3}{-2}=\dfrac{z}{1}$}
    {\True Giả sử sau $t$ (s) kể từ lúc xuất phát $(t \geq 0)$, cabin đến điểm $M$. Khi đó điểm $M$ có tọạ độ là $\left(3 t+10 ;-3 t+3 ;\dfrac{3t}{2}\right)$}
    {Cabin dừng ở điểm $B$ có hoành độ $x_B=550$. Khi đó $AB=750$ m}
    {\True Đường cáp $AB$ tạo với mặt phẳng $(Oxy)$ góc $19^\circ$ (kết quả làm tròn đến hàng đơn vị của độ)}
    \loigiai{
        \begin{itemchoice}
            \itemch Phương trình chính tắc của đường cáp là $\dfrac{x-10}{2}=\dfrac{y-3}{-2}=\dfrac{z}{1}$.
            \itemch Do tốc độ chuyển động của cabin là $4{,}5$ m/s nên độ dài $A M$ bằng $4{,}5 t$ m. \\
            Vì vậy $\left|\vec{AM}\right|=4{,}5 t$ $(t \geq 0)$.\\
            Do hai véc-tơ $\vec{A M}$ và $\vec{u}$ là cùng phương và cùng hướng nên $\vec{A M}=k \vec{u}$ với $k$ là số thực dương nào đó. \\
            Suy ra $\left|\vec{A M}\right|=k|\vec{u}|=k \cdot \sqrt{2^2+(-2)^2+1}=3 k$. Do đó $3 k=4{,}5 t$. Suy ra $k=\dfrac{3 t}{2}$. \\
            Vì thế, ta có $\vec{A M}=\dfrac{3 t}{2} \vec{u}=\left(3 t ;-3 t ; \dfrac{3 t}{2}\right)$.\\
            Gọi tọa độ của điểm $M$ là $\left(x_M ; y_M ; z_M\right)$.\\
            Ta có $\vec{A M}=\left(x_M-x_A ; y_M-y_A ; z_M-z_A\right)=\left(3 t ;-3 t ; \dfrac{3 t}{2}\right)$.\\
            Nên $\heva{&x_M=3 t+x_A \\ &y_M=-3 t+y_A \\ &z_M=\dfrac{3 t}{2}+z_A}\Leftrightarrow\heva{&x_M=3 t+10 \\ &y_M=-3 t+3 \\ &z_M=\dfrac{3 t}{2}.}$\\
            Vậy điểm $M$ có tọạ độ là $\left(3 t+10 ;-3 t+3 ;\dfrac{3 t}{2}\right)$.
            \itemch Do $x_B=550$ nên $3 t+10=550$, tức là $t=180$ s. Do đó, ta có điểm $B(550 ;-537 ; 270)$. \\
            Vậy $A B=\sqrt{(550-10)^2+(-537-3)^2+(270-0)^2}=\sqrt{656100}=810$ m.
            \itemch Đường thẳng $AB$ có véc-tơ chỉ phương $\vec{u}=(2 ;-2 ; 1)$ và mặt phẳng $(Oxy)$ có véc-tơ pháp tuyến $\vec{k}=(0 ; 0 ; 1)$. Do đó, ta có
            $$
            \sin (\Delta,(O x y))=|\cos (\vec{u},\vec{k})|=\dfrac{\left|\vec{u} \cdot \vec{k}\right|}{|\vec{u}| \cdot|\vec{k}|}=\dfrac{1}{3 \cdot 1}=\dfrac{1}{3}.$$
            Vậy $(\Delta,(O x y)) \approx 19^{\circ}$.
        \end{itemchoice}
    }
\end{ex}

\begin{ex}%Câu 15
Một đoàn tàu đang đứng yên trong sân ga, ngay trước đầu tàu có một cái cây. Đoàn tàu khởi hành từ trạng thái đứng yên với gia tốc $ a=0,005t\text(\text{m/}{\text{s}^{\text{2}}})$ và đi qua cái cây trong thời gian $60$ giây. Sau $80$ giây đoàn tàu chuyển sang trạng thái chuyển động đều. Xét tính đúng sai của các khẳng định sau:\\
\centerline{\includegraphics[width=.6\textwidth]{images/4.15}}
\choiceTF
{Vận tốc của đoàn tàu là $v=5.10^{-3}{t^2}\text(\text{m/s})$}
{\True Chiều dài của đoàn tàu là $l=180$m}
{\True Sau $80$ giây, đoàn tàu chuyển động với tốc tốc $57,6\text(\text{km/h})$}
{Sau khi chuyển động đều một thời gian, đoàn tàu gặp một cây cầu có chiều dài $480$ m, khi đó đoàn tàu đi qua cây cầu đó trong thời gian $30$ giây}
\loigiai{

a) Sai: Vận tốc của tàu là $ v(t)=\displaystyle\int{a(t)\mathrm{\,d}t}=\displaystyle\int{0,005t\,\mathrm{\,d}t}=2,5.10^{-3}{t^2}\,(\text{m/s})$ với $ v(0)=0$\\
b) Đúng: Chiều dài của đoàn tàu bằng quãng đường tàu đi trong 60 giây đầu tiên và bằng\\
$ l=\displaystyle\int\limits_0^{60}{v(t)\mathrm{\,d}t}=\displaystyle\int\limits_0^{60}{2,5.10^{-3}{t^2}\mathrm{\,d}t}=180$ m\\
c) Đúng: Sau 80 giây, đoàn tàu chuyển động với tốc tốc:\\
$ v\left(80\right)=2,5.10^{-3}{80^2}=16\,\,\left(\text{m}/\text{s}\right)=57,6\left(\text{km}/\text{h}\right)$\\
d) Sai: Sau khi chuyển động đều một thời gian, đoàn tàu gặp một cây cầu có chiều dài 480 (m).\\
Khi đó đoàn tàu đi qua cây cầu đó trong thời gian $ t=\dfrac{480+180}{16}=41,25$ giây.}
\end{ex}

\begin{ex}%Câu 16
Có hai hộp đựng phiếu thi, mỗi phiếu ghi một câu hỏi. Hộp thứ nhất có $15$ phiếu và hộp thứ hai có $9$ phiếu. Sinh viên A đi thi chỉ thuộc $10$ câu ở hộp thứ nhất và $8$ câu ở hộp thứ hai (Kết quả làm tròn đến 2 chữ số sau dấu phẩy)
\choiceTF
{\True Thầy giáo rút ngẫu nhiên từ mỗi hộp ra một phiếu thi, sau đó cho sinh viên A rút ngẫu nhiên ra 1 phiếu từ 2 phiếu mà thầy giáo đã rút. Xác suất để sinh viên A trả lời được câu hỏi trong phiếu là $0,78$}
{Thầy giáo rút ngẫu nhiên ra 1 phiếu từ hộp thứ nhất bỏ vào hộp thứ hai, sau đó cho sinh viên A rút ngẫu nhiên ra 1 phiếu từ hộp thứ hai. Xác suất để sinh viên trả lời được câu hỏi trong phiếu là $0,73$}
{Thầy giáo rút ngẫu nhiên ra 2 phiếu từ hộp thứ nhất bỏ vào hộp thứ hai, sau đó cho sinh viên A rút ngẫu nhiên ra 2 phiếu từ hộp thứ hai, xác suất để sinh viên đó rút được hai câu thuộc là $0,62$}
{Thầy giáo rút ngẫu nhiên ra 1 phiếu từ hộp thứ nhất bỏ vào hộp thứ hai, sau đó cho sinh viên A rút ngẫu nhiên ra 2 phiếu từ hộp thứ hai, xác suất để sinh viên đó rút được hai câu thuộc là $0,83$}
\loigiai{
a) Đúng: Gọi $A_1$ là biến cố sinh viên rút ra được phiếu thuộc hộp 1.\\
$A_2$ là biến cố sinh viên rút ra được phiếu thuộc hộp 2.\\
$A$ là biến cố sinh viên rút ra $1$ câu thuộc, khi đó:\\
$A=\left(A_1\cap A\right)\cup\left(A_2\cap A\right)\Rightarrow P(A)=P\left(A_1\right).P\left(A|A_1\right)+P\left(A_2\right).P\left(A|A_2\right)$ .\\
Ta có: $P\left(A_1\right)=\dfrac{1}{2};P\left(A_2\right)=\dfrac{1}{2}\Rightarrow P\left(A|A_1\right)=\dfrac{C_{10}^1}{C_{15}^1}=\dfrac{2}{3};P\left(A|A_2\right)=\dfrac{C_8^1}{C_9^1}=\dfrac{8}{9}$
Vậy $P(A)=\dfrac{7}{9}\approx 0,78$ .\\
b) Sai: Gọi $B_1$ là biến cố thầy giáo rút 1 câu thuộc từ hộp 1 bỏ vào hộp 2, khi đó hộp 2 có 9 câu thuộc và 1 câu không thuộc.\\
Gọi $B_2$ là biến cố thầy giáo rút 1 câu không thuộc từ hộp 1 bỏ vào hộp 2, khi đó hộp 2 có 8 câu thuộc và 2 câu không thuộc.\\
Gọi $B$ là biến cố sinh viên rút ra $1$ câu thuộc, khi đó:\\
$B=\left(B_1\cap B\right)\cup\left(B_2\cap B\right)\Rightarrow P(B)=P\left(B_1\right).P\left(B|B_1\right)+P\left(B_2\right).P\left(B|B_2\right)$ .\\
Ta có: $P\left(B_1\right)=\dfrac{C_{10}^1}{C_{15}^1}=\dfrac{2}{3};P\left(B_2\right)=\dfrac{C_5^1}{C_{15}^1}=\dfrac{1}{3}\Rightarrow P\left(B|B_1\right)=\dfrac{C_9^1}{C_{10}^1}=\dfrac{9}{10};P\left(B|B_2\right)=\dfrac{C_8^1}{C_{10}^1}=\dfrac{4}{5}$
Vậy $P(B)\approx 0,94$ .\\
c) Sai: Gọi $C_1$ là biến cố thầy giáo rút 2 câu thuộc từ hộp 1 bỏ vào hộp 2, khi đó hộp 2 có 10 câu thuộc và 1 câu không thuộc.\\
Gọi $C_2$ là biến cố thầy giáo rút 1 câu thuộc và 1 câu không thuộc từ hộp 1 bỏ vào hộp 2, khi đó hộp 2 có 9 câu thuộc và 2 câu không thuộc.\\
Gọi $C_3$ là biến cố thầy giáo rút 2 câu không thuộc từ hộp 1 bỏ vào hộp 2, khi đó hộp 2 có 8 câu thuộc và 3 câu không thuộc.\\
Gọi $C$ là biến cố sinh viên rút ra $2$ câu thuộc, khi đó: $C=\left(C_1\cap C\right)\cup\left(C_2\cap C\right)\cup\left(C_3\cap C\right)$\\
$\Rightarrow P(C)=P\left(C_1\right).P\left(C|C_1\right)+P\left(C_2\right).P\left(C|C_2\right)+P\left(C_3\right).P\left(C|C_3\right)$\\
Ta có: $P\left(C_1\right)=\dfrac{C_{10}^2}{C_{15}^2}=\dfrac{3}{7};P\left(C_2\right)=\dfrac{C_5^1.C_{10}^1}{C_{15}^2}=\dfrac{10}{21};P\left(C_3\right)=\dfrac{C_5^2}{C_{15}^2}=\dfrac{2}{21}$\\
$P\left(C|C_1\right)=\dfrac{C_{10}^2}{C_{11}^2}=\dfrac{9}{11};P\left(C|C_2\right)=\dfrac{C_9^2}{C_{11}^2}=\dfrac{12}{35};P\left(C|C_3\right)=\dfrac{C_8^2}{C_{11}^2}=\dfrac{3}{55}$
Vậy $P(C)\approx 0,52$ .\\
d) Sai: Gọi $D_1$ là biến cố thầy giáo rút 1 câu thuộc từ hộp 1 bỏ vào hộp 2, khi đó hộp 2 có 9 câu thuộc và 1 câu không thuộc.\\
Gọi $D_2$ là biến cố thầy giáo rút 1 câu không thuộc từ hộp 1 bỏ vào hộp 2, khi đó hộp 2 có 8 câu thuộc và 2 câu không thuộc.\\
Gọi $D$ là biến cố sinh viên rút ra 2 câu thuộc, khi đó:\\
$D=\left(D_1\cap D\right)\cup\left(D_2\cap D\right)\Rightarrow P(D)=P\left(D_1\right).P\left(D|D_1\right)+P\left(D_2\right).P\left(D|D_2\right)$ .\\
Ta có: $P\left(D_1\right)=\dfrac{C_{10}^1}{C_{15}^1}=\dfrac{2}{3};P\left(D_2\right)=\dfrac{C_5^1}{C_{15}^1}=\dfrac{1}{3}$ $\Rightarrow P\left(D|D_1\right)=\dfrac{C_9^2}{C_{10}^2}=\dfrac{4}{5};P\left(D|D_2\right)=\dfrac{C_8^2}{C_{10}^2}=\dfrac{28}{45}$ .\\
Vậy $P(D)\approx 0,74$ }
\end{ex}
\Closesolutionfile{ans}
\TNSA
\Opensolutionfile{ans}[ans/ans-4-SA]
\begin{ex}%Câu 17
Cho hình lăng trụ đứng $ ABC.A'B'C'$ có đáy $ ABC$ là tam giác vuông tại $ B$ và có độ dài các cạnh$ AB=a\sqrt{3}$, $ BC=2a$, $ AA'=a\sqrt{2}$. Gọi $ M$ là trung điểm của $ BC$. Tính khoảng cách giữa hai đường thẳng $ AM$ và $B'C$ khi $ a=1$ (Kết quả làm tròn đến hàng phần trăm).
\shortans{0,55}
\loigiai{
{\color{red}HÌNH Ở ĐÂY}\\
Gọi $ N$ là trung điểm của $ B{B}'$ thì $ MN\text{//}{B}'C\Rightarrow{B}'C\text{//}\left(AMN\right)$.\\
Ta có: $ d\left(B'C,AM\right)=d\left(B'C,\left(AMN\right)\right)=d\left(B',\left(AMN\right)\right)=d\left(B,\left(AMN\right)\right)$.\\
Dựng $ BI\perp AM, BH\perp NI$$\Rightarrow BH\perp\left(AMN\right)$ nên do đó $ d\left(B'C,AM\right)=d\left(B,\left(AMN\right)\right)=BH$.\\
Vì $\Delta ABM$ vuông tại $ B$, ta có $BM=\dfrac{BC}{2}=a$ ; $ BI=\dfrac{BA.BM}{\sqrt{B{A^2}+B{M^2}}}=\dfrac{a\sqrt{3}.a}{\sqrt{3a^2+a^2}}=\dfrac{a\sqrt{3}}{2}$.\\
Xét $\Delta BIN$ vuông tại $ B$, ta có:\\
$ BN=\dfrac{B{B}'}{2}=\dfrac{a\sqrt{2}}{2}$; $ BH=\dfrac{BN.BI}{\sqrt{B{N^2}+B{I^2}}}=\dfrac{\dfrac{a\sqrt{2}}{2}.\dfrac{a\sqrt{3}}{2}}{\sqrt{\left(\dfrac{a\sqrt{2}}{2}\right)^2+\left(\dfrac{a\sqrt{3}}{2}\right)^2}}=\dfrac{a\sqrt{30}}{10}$.\\
Vậy $ d\left(B'C,AM\right)=BH=\dfrac{\sqrt{30}}{10}\xrightarrow{a=1}d\left(B'C,AM\right)\approx 0,55$.}
\end{ex}

\begin{ex}%Câu 18
Trong một cuộc thi về “bữa ăn dinh dưỡng”, ban tổ chức yêu cầu để đảm bảo lượng dinh dưỡng hằng ngày thì mỗi gia đình có $4$ thành viên cần ít nhất $900$ đơn vị prôtêin và $400$ đơn vị lipít trong thức ăn hằng ngày. Mỗi kg thịt bò chứa $800$ đơn vị prôtêin và $200$ đơn vị lipit, $1\text{(kg)}$ thịt heo chứa $600$ đơn vị prôtêin và $400$ đơn vị lipit. Biết rằng người nội trợ chỉ được chi tối đa $200$ ngàn đồng để mua thịt. Biết rằng $1\text{(kg)}$ thịt bò giá $200$ ngàn đồng, $1\text{(kg)}$ thịt heo giá $100$ ngàn đồng. Người nội trợ nên mua $x\text{(kg)}$ thịt bò và $y\text{(kg)}$ thịt heo để phí thấp nhất cho khẩu phần thức ăn mà vẫn đảm bảo chất dinh dưỡng. Khi đó hãy tìm $x+2y.$
\shortans{3}
\loigiai{
Chi phí mua thịt là $ F\left(x;y\right)=200x+100y$ (ngàn đồng).\\
Hệ điều kiện ràng buộc giữa $ x$ và $ y$ là $\left\{\begin{matrix}
x\ge 0\\
y\ge 0\\
800x+600y\ge 900\\
200x+400y\ge 400\\
\end{matrix}\Leftrightarrow\left\{\begin{matrix}
x\ge 0\\
y\ge 0\\
8x+6y\ge 9\\
x+2y\ge 2\\
\end{matrix}\right.\right.$.\\
Miền nghiệm của hệ bất phương trình này là miền không bị gạch sọc trong hình vẽ dưới đây\\
{\color{red}HÌNH Ở ĐÂY}\\
Các điểm cực biên là $ A\left(0;\,1,5\right),\,\,B\left(0,6;\,0,7\right),\,\,C\left(2;\,0\right)$.\\
Ta có: $ F\left(0;1,5\right)=150,\,\,\,F\left(0,6;0,7\right)=190,\,\,\,F\left(2;0\right)=400$.\\
Vậy chi phí mua thịt thấp nhất khi $ x=0,y=1,5\Rightarrow x+2y=3$.
}
\end{ex}
%
\begin{ex}%Câu 19
Đường đi của một khinh khí cầu được gắn trong hệ trục tọa độ là một đường cong bậc hai trên bậc nhất có đồ thị cắt trục hoành tại hai điểm có tọa độ là $\left(1;0\right)$ và $\left(8;0\right)$ với đơn vị trên hệ trục tọa độ là $ 1$ km. Biết rằng điểm cực đại của đồ thị hàm số là điểm $\left(6;5\right).$ Hỏi khi khí cầu đi qua điểm cực đại và cách mặt đất $3875$ M thì khí cầu cách gốc tọa độ theo phương ngang bao nhiêu? (đơn vị: km)\\
\centerline{\includegraphics[width=.3\textheight]{images/4.19}}
\shortans{7,2}
\loigiai{
Đồ thị hàm số bậc hai trên bậc nhất $ y=\dfrac{a{x^2}+bx+c}{x+m}$ cắt trục hoành tại hai điểm có hoành độ $ x=1,\,\,x=8$ nên ta có: $ y=\dfrac{a\left(x-1\right)\left(x-8\right)}{x+m}=\dfrac{a\left(x^2-9x+8\right)}{x+m}.$\\
Đường thẳng qua hai điểm cực trị của đồ thị hàm số là $(d):y=\dfrac{\left[a\left(x^2-9x+8\right)\right]'}{\left(x+m\right)'}=a\left(2x-9\right)$ qua điểm cực đại $\left(6;5\right)\Leftrightarrow 5=a\left(2.6-9\right)\Leftrightarrow a=\dfrac{5}{3}\Rightarrow y=\dfrac{5\left(x^2-9x+8\right)}{3\left(x+m\right)}$.\\
Ta có $y'(6)=0\Leftrightarrow m=-\dfrac{28}{3}\Rightarrow y=\dfrac{5\left(x^2-9x+8\right)}{3\left(x-\dfrac{28}{3}\right)}$.\\
Xét $ y=3,875\Leftrightarrow\dfrac{5\left(x^2-9x+8\right)}{3\left(x-\dfrac{28}{3}\right)}=3,875\Leftrightarrow\left[\begin{aligned}
& x\approx 4,125\\ 
& x\approx 7,2\\ 
\end{aligned}\right.$\\
Vậy khi khí cầu đi qua điểm cực đại và cách mặt đất $ 3875$(m) thì khí cầu cách gốc tọa độ theo phương ngang $ 7,2$ km.}
\end{ex}

\begin{ex}%Câu 20
Hệ thống lọc nước bể bơi vô cùng quan trọng khi tiến hành xây dựng công trình bơi lội để nguồn nước được làm sạch thường xuyên và giữ vệ sinh cho người bơi. Trong quá trình vận hành lọc nước thì lượng nước trong bể sẽ thay đổi theo thời gian. Lượng nước trong bể giảm nếu hệ thống đang xả nước bẩn ra khỏi bể và tăng nếu hệ thống đang cấp thêm nước sạch cho bể. Biết rằng $1$ gallon gần bằng $3,785$ lít, dung tích của bể là $ 1000$ gallon và thời điểm $6$ giờ sáng bể chứa $ 250$ gallon nước. Hàm số $ f(t)$ biểu thị cho tốc độ thay đổi lượng nước trong bể theo thời gian $ t$ giờ, từ thời điểm $6$ giờ sáng đến $6$ giờ chiều được cho bởi $f(t)=\left\{\begin{aligned}
& 100t, \text{ khi }0\le t\le 3\\ 
& 900-200t, \text{ khi }3\le t\le 6\\ 
& 100t-900, \text{ khi }6\le t\le 12\\ 
\end{aligned}\right.$ với mốc thời gian $ t=0$ tại thời điểm $6$ giờ sáng. Hỏi ở thời điểm $6$ giờ chiều thì trong bể chứa nhiêu gallon nước?
\shortans{700}
\loigiai{
Gọi $ F(t)$ là một nguyên hàm của hàm số $ f(t)$\\
Vì $ f(t)$ biểu thị cho tốc độ thay đổi lượng nước trong bể theo thời gian $ t$ nên $ F(t)$ chính là lượng nước có trong bể theo thời gian $ t$.\\
Lượng nước trong bể lúc 6 giờ sáng (ứng với $ t=0$) là $ F(0)=250$.\\
Lượng nước trong bể lúc 6 giờ chiều (ứng với $ t=12$) là $ F\left(12\right)=F(0)+\displaystyle\int\limits_0^{12}{f(t)\mathrm{\,d}t}$\\
$=250+\displaystyle\int\limits_0^3100t\mathrm{\,d}t+\displaystyle\int\limits_3^6\left(900-200t\right)\mathrm{\,d}t+\displaystyle\int\limits_6^{12}{\left(100t-900\right)\mathrm{\,d}t}=700$ (galon)\\
Đáp án:\\
7 0 0}
\end{ex}

\begin{ex}%Câu 21
Hệ thống định vị toàn cầu GPS (Global Positioning System) là một hệ thống cho phép xác định vị trí của một vật thể trong không gian. Trong cùng một thời điểm vị trí của một điểm $ M$ trong không gian sẽ được xác định bởi bốn vệ tinh cho trước nhờ các bộ thu phát tín hiệu đặt trên các vệ tinh đó. Giả sử trong không gian với hệ trục tọa độ $ Oxyz$, có bốn vệ tinh lần lượt đặt tại các điểm $ A\left(2;4;0\right),\,B\left(0;4;6\right),\,C\left(2;0;6\right),\,D\left(-1;-2;-3\right)$ và vị trí của điểm $ M\left(a;b;c\right)$ thỏa mãn biểu thức $ MA+MB+MC+MD$ nhỏ nhất. Tính độ dài $ MO$ (kết quả làm tròn đến hàng phần chục)
\shortans{3,7}
\loigiai{
Gọi $ G$ là trọng tâm của tứ diện $OABC$. Suy ra $G(1;2;3)$.\\
Từ đó suy ra $ GA=GB=GC=GO=\sqrt{14}$ và $O$ là trung điểm $GD$.
\begin{align*}
MA+MB+MC+MD&=\dfrac{MA \cdot GA + MB \cdot GB + MC \cdot GC + MD \cdot GD}{GA} \\
& \ge \dfrac{\vec{MA} \cdot \vec{GA}+\vec{MB} \cdot \vec{GB}+\vec{MC} \cdot \vec{GC}+\vec{MD} \cdot \vec{GO}}{GA} \\
&=\dfrac{\vec{MG}\left(\vec{GA}+\vec{GB}+\vec{GC}+\vec{GO}\right)+5GA^2}{GA}=5GA
\end{align*}
Dấu bằng xảy ra khi và chỉ khi $ M$ trùng với điểm $ G$.\\
Khi đó $MO=GO=\sqrt{14}\approx 3,7 $.
}
\end{ex}

\begin{ex}%Câu 22
Một hộp chứa $10$ viên bi xanh và $5$ viên bi đỏ. Bạn An lấy ra ngẫu nhiên $1$ viên bi từ hộp, xem màu, rồi bỏ ra ngoài. Nếu viên bi An lấy ra có màu xanh, bạn Bình sẽ lấy ra ngẫu nhiên $2$ viên bi từ hộp; còn nếu viên bi An lấy ra có màu đỏ, bạn Bình sẽ lấy ra ngẫu nhiên $3$ viên bi từ hộp. Tính xác suất để An lấy được viên bi màu xanh, biết rằng tất cả các viên bi được hai bạn chọn ra đều có đủ cả hai màu.
\shortans{0,55}
\loigiai{
Gọi $ A$ là biến cố An lấy được viên bi màu xanh thì $\overline{A}$ là biến cố An lấy được viên bi màu đỏ.
Gọi $ B$ là biến cố tất cả các viên bi được hai bạn chọn ra đều có đủ cả hai màu.\\
Ta có $ P(A)=\dfrac{10}{15},P\left(\overline{A}\right)=\dfrac{5}{15}$. Ta cần tính $ P\left(A\mid B\right)=\dfrac{P\left(AB\right)}{P(B)}$.\\
Khi $ A$ xảy ra thì trong hộp còn 9 viên bi xanh và 5 viên bi đỏ, Bình cần lấy ra 2 viên bi trong đó có ít nhất một viên bi đỏ nên $ P\left(B\mid A\right)=\dfrac{C_{14}^2-C_9^2}{C_{14}^2}$.\\
Khi $\overline{A}$ xảy ra thì trong hộp còn 10 viên bi xanh và 4 viên bi đỏ, Bình cần lấy ra 3 viên bi trong đó có ít nhất một viên bi xanh nên $ P\left(B\mid\overline{A}\right)=\dfrac{C_{14}^3-C_4^3}{C_{14}^3}$.\\
Vậy $ P\left(AB\right)=P(A).P\left(B\mid A\right)=\dfrac{10}{15}\cdot\dfrac{\left(C_{14}^2-C_9^2\right)}{C_{14}^2}=\dfrac{110}{273}$\\
Mặt khác: $ P(B)=P(A).P\left(B\mid A\right)+P\left(\overline{A}\right).P\left(B\mid\overline{A}\right)$ $=\dfrac{10}{15}\cdot\dfrac{\left(C_{14}^2-C_9^2\right)}{C_{14}^2}+\dfrac{5}{15}\cdot\dfrac{\left(C_{14}^3-C_4^3\right)}{C_{14}^3}=\dfrac{200}{273}$.\\
Vậy $ P\left(A\mid B\right)=\dfrac{\dfrac{110}{273}}{\dfrac{200}{273}}=0,55$.}
\end{ex}
\Closesolutionfile{ans}
\Closesolutionfile{ansbook}
\inputansbox{6,2,3}{ans/ans-4-T,ans/ans-4-TF,ans/ans-4-SA}